% sections/07_falsifiability.tex
% JPHQ-04 — Falsifiability conditions (OC/ETS-internal)

\subsection*{Falsifiability Conditions}

The No-Go theorems presented in this preprint are \emph{falsifiable} in a strict
and formal sense. Each impossibility claim depends on explicitly stated
assumptions within the \OC/\kw{ETS.K\_EA.v1.1} framework. Failure of any such
assumption invalidates the corresponding No-Go result.

All falsifiability conditions are formulated entirely at the OC/ETS level.
No appeal is made to empirical data, managerial interpretation, behavioral
assumptions, or external semantic layers.

\paragraph{F1. Axis Embedding without Incompatibility.}
A No-Go theorem based on axis incompatibility is falsified if there exists an
OC-admissible embedding
\[
\varphi:\; A_{\mathrm{new}} \hookrightarrow \Ax(\KEA)
\]
such that the extended axis set remains compatible with all thresholds
$\Th(\KEA)$ and admits at least one admissible state in $\Om(\KEA)$.
Existence of such an embedding refutes claims of unavoidable structural
axis incompatibility.

\paragraph{F2. Threshold Reconfiguration with Non-Empty State Space.}
A No-Go theorem relying on a threshold viability barrier is falsified if there
exists a threshold reconfiguration $\Th'(\KEA)$ such that
\[
\Om(\KEA;\Th') \neq \varnothing,
\]
i.e., the admissible state set remains non-empty after the threshold update.
Such a configuration demonstrates that the assumed threshold barrier is not
structurally necessary.

\paragraph{F3. Cycle Preservation across Structural Change.}
A No-Go theorem based on cycle-closure obstruction is falsified if there exists
a family of admissible cycles $\Cyc'(\KEA)$ compatible with the transformed
structural type and satisfying
\[
\kcont(\KEA) > 0
\]
throughout the postulated transformation.
Existence of such cycles refutes claims that transformation necessarily entails
loss of continuity or forced entry into phase $\PHST$.

\paragraph{F4. Admissible Phase-Transition Path.}
Phase-incompatibility No-Go results are falsified if there exists an explicit,
OC/ETS-admissible phase-transition path from the initial phase of $\KEA$ to a
phase compatible with the transformed structural type, without traversing any
forbidden transition.
Existence of such a path invalidates the assumed phase obstruction.

\paragraph{F5. Boundary-Compatible Trajectories.}
Boundary-based No-Go theorems are falsified if there exists an admissible
trajectory that crosses $\dOm(\KEA)$ in a manner permitted by OC/ETS and reaches
a structurally transformed state without violating admissibility.
This demonstrates that the boundary constraints do not constitute an absolute
structural obstruction.

\begin{remark}
Falsification of any No-Go theorem therefore requires an explicit
\emph{OC/ETS-internal counterexample}: a demonstrable axis embedding, threshold
configuration, admissible cycle family, phase-transition path, or
boundary-respecting trajectory. In the absence of such constructions, the
corresponding impossibility results remain valid within the assumed formal
system.
\end{remark}
